% =====================================================================
%  Thème 2 — Angles associés — NIVEAU FACILE
%  Corrigé
% =====================================================================
\documentclass[11pt,a4paper]{article}
\usepackage[utf8]{inputenc}
\usepackage[T1]{fontenc}
\usepackage{lmodern}
\usepackage[french]{babel}
\usepackage{amsmath,amssymb,mathtools}
\usepackage{icomma}
\usepackage{xcolor}
\usepackage{colortbl}
\usepackage{tikz}
\usepackage[most]{tcolorbox}
\usepackage{enumitem}
\usepackage{array}
\usepackage[a4paper,margin=2.2cm]{geometry}
\usetikzlibrary{arrows.meta,calc}

\definecolor{primary}{RGB}{214,51,108}
\definecolor{primarydark}{RGB}{140,20,64}
\definecolor{excol}{RGB}{0,135,90}

\DeclareMathOperator{\tg}{tg}
\DeclareMathOperator{\cotg}{cotg}
\newcommand{\degr}{^{\circ}}
\newcommand{\Z}{\mathbb{Z}}

\newcounter{exo}
\newcommand{\exo}[1]{\medskip\refstepcounter{exo}%
  \noindent{\large\bfseries\color{primarydark}Exercice \theexo}%
  \ \textcolor{primary}{\textbullet}\ \textbf{#1}\par\nopagebreak\smallskip}
\newcommand{\rep}{\textcolor{excol}{$\blacktriangleright$}\ }

\setlength{\parindent}{0pt}
\pagestyle{empty}

\begin{document}

\begin{center}
{\LARGE\bfseries\color{primary}Thème 2 — Angles associés}\\[2pt]
{\color{primarydark}\textsc{Niveau Facile — Corrigé détaillé}}
\end{center}
\hrule height 1pt
\medskip

\exo{Opposés ($-\alpha$)}
Symétrie par l'axe des cosinus : $\cos$ est \emph{pair}, $\sin$ et $\tg$ sont \emph{impaires}.
\begin{itemize}[leftmargin=1.4em,itemsep=1pt]
\item[\textbf{(a)}] \rep $\cos(-\alpha)=\cos\alpha$.
\item[\textbf{(b)}] \rep $\sin(-\alpha)=-\sin\alpha$.
\item[\textbf{(c)}] \rep $\tg(-\alpha)=-\tg\alpha$.
\item[\textbf{(d)}] \rep $\sin(-\alpha)+\sin\alpha=-\sin\alpha+\sin\alpha=0$.
\end{itemize}

\exo{Supplémentaires ($\pi-\alpha$)}
Symétrie par l'axe des sinus : même $\sin$, $\cos$ et $\tg$ changent de signe.
\begin{itemize}[leftmargin=1.4em,itemsep=1pt]
\item[\textbf{(a)}] \rep $\sin(\pi-\alpha)=\sin\alpha$.
\item[\textbf{(b)}] \rep $\cos(\pi-\alpha)=-\cos\alpha$.
\item[\textbf{(c)}] \rep $\tg(\pi-\alpha)=-\tg\alpha$.
\item[\textbf{(d)}] \rep $\sin\!\left(\pi-\dfrac{\pi}{6}\right)=\sin\dfrac{\pi}{6}=\dfrac12$
(deux angles supplémentaires ont le même sinus).
\end{itemize}

\exo{Anti-supplémentaires ($\pi+\alpha$)}
Symétrie centrale (par $O$) : $\cos$ et $\sin$ changent de signe, la $\tg$ est inchangée (période $\pi$).
\begin{itemize}[leftmargin=1.4em,itemsep=1pt]
\item[\textbf{(a)}] \rep $\cos(\pi+\alpha)=-\cos\alpha$.
\item[\textbf{(b)}] \rep $\sin(\pi+\alpha)=-\sin\alpha$.
\item[\textbf{(c)}] \rep $\tg(\pi+\alpha)=\tg\alpha$.
\item[\textbf{(d)}] \rep $\tg\!\left(\pi+\dfrac{\pi}{4}\right)=\tg\dfrac{\pi}{4}=1$.
\end{itemize}

\exo{Complémentaires et anti-complémentaires ($\frac{\pi}{2}\pm\alpha$)}
Autour de $\frac{\pi}{2}$, les fonctions $\sin$ et $\cos$ \textbf{s'échangent}.
\begin{itemize}[leftmargin=1.4em,itemsep=1pt]
\item[\textbf{(a)}] \rep $\sin\!\left(\dfrac{\pi}{2}-\alpha\right)=\cos\alpha$ \emph{(complémentaires)}.
\item[\textbf{(b)}] \rep $\cos\!\left(\dfrac{\pi}{2}-\alpha\right)=\sin\alpha$ \emph{(complémentaires)}.
\item[\textbf{(c)}] \rep $\cos\!\left(\dfrac{\pi}{2}+\alpha\right)=-\sin\alpha$ \emph{(anti-complémentaires)}.
\item[\textbf{(d)}] \rep $\sin\!\left(\dfrac{\pi}{2}+\alpha\right)=\cos\alpha$ \emph{(anti-complémentaires)}.
\end{itemize}

\exo{Reconnaître la nature de deux angles}
On compare la somme (ou la différence) des deux angles à $90\degr$, $180\degr$, $0\degr$.
\begin{itemize}[leftmargin=1.4em,itemsep=1pt]
\item[\textbf{(a)}] \rep $30\degr+150\degr=180\degr$ : \textbf{supplémentaires}.
\item[\textbf{(b)}] \rep $40\degr+50\degr=90\degr$ : \textbf{complémentaires}.
\item[\textbf{(c)}] \rep $70\degr$ et $-70\degr$ sont \textbf{opposés} (somme nulle).
\item[\textbf{(d)}] \rep $200\degr=180\degr+20\degr$ : \textbf{anti-supplémentaires}.
\end{itemize}

\end{document}
